
\section{14.05.2026 - Gute argumente kann man immer gebrauchen}

\subsection{Wortschatz}

\begin{vocabtable}
\deword{der Bezug, die Bezüge}{reference / connection / relation}{In meiner Antwort stelle ich einen klaren Bezug zur Stelle her.}
\deword{das Redemittel, die Redemittel}{useful phrase / expression}{Für eine gute Argumentation helfen passende Redemittel.}
\deword{die Verhandlung, die Verhandlungen}{negotiation}{Gute Argumente sind in Verhandlungen besonders wichtig.}
\deword{zitieren + akk}{to quote / to cite}{In einem Vortrag kann man eine Studie zitieren.}
\deword{veraltet}{outdated / obsolete}{Diese Information ist veraltet und sollte überprüft werden.}
\deword{sich erkundigen nach + dat}{to inquire about / to ask about}{Ich erkundige mich nach den Anforderungen der Stelle.}
\deword{eineinhalb}{one and a half}{Das Projekt hat eineinhalb Monate gedauert.}
\deword{einreichen + akk}{to submit}{Ich habe meine Bewerbung gestern eingereicht.}
\deword{rechnen mit + dat}{to expect / to reckon with}{Wir müssen mit schwierigen Fragen rechnen.}
\end{vocabtable}
